% Einheit 4 von 4 – Scharen am Körper (bank/scharen-von-geraden-und-ebenen/e4.jsonl, zusammenbau v0.1)
%% TODO Zeitmarke Einheit 4: Marken-Zeile nicht lesbar
%% TODO Zweigzeile Teil 1: was hier gelernt wird (Satz in Schülersprache aus dem Einheitstitel) – Einheit 4
\einheitenkopf[e4]{Einheit 4 von 4 · Scharen am Körper}
\zweigzeile{Abitur LK}
\setcounter{aufgabe}{15}

%% TODO Titel: Ich-kann-Satz fehlt in der Bank (Platzhalter: Kettenname) – Nr. 16
%% TODO Anweisung: ein Satz über der ersten Teilaufgabe fehlt in der Bank – Nr. 16
\begin{aufgabe}{Scharen am Körper}
%% TODO \rechenplatz{n}: Schrittzahl des Grundfalls fehlt in der Bank (nur bei fester Schreibform, 2.3 b)
\begin{teile}
\teil Die Abbildung zeigt den Würfel mit $0 \le x_1, x_2, x_3 \le 3$. Bei welchen Werten läuft $E_k: x_1 + x_2 + x_3 = k$ durch Ecken? \\ \kreuz{0, 3, 6 und 9} \\ \kreuz{nur 0 und 9} \\ \kreuz{0, 1, 2 und 3}

\begin{ksys3}\rquader{0,0,0}{3}{3}{3}\end{ksys3}
\teil Die Abbildung zeigt den Würfel mit $0 \le x_1, x_2, x_3 \le 4$; $E_k: x_1 + x_2 + x_3 = k$. Gib $E_2$ an und zeichne die Schnittfigur von $E_2$ mit dem Würfel ein.

\begin{ksys3}\rquader{0,0,0}{4}{4}{4}\end{ksys3}
\teil Die Ebene $E: x_1 + x_2 + x_3 = 10$ schneidet den Würfel mit $0 \le x_1, x_2, x_3 \le 4$ in einem Dreieck mit den Ecken $P(4 | 4 | 2)$ und $Q(4 | 2 | 4)$. Bestimme die dritte Ecke $R$. $R($ \leerfeld | \leerfeld | \leerfeld $)$
\teil Gib für $0 < k < 6$ an, wie viele Ecken die Schnittfigur von $E_k: x_1 + x_2 + x_3 = k$ mit dem Würfel $0 \le x_1, x_2, x_3 \le 2$ hat.
\teil Die Ebenen $E_k: x_1 + x_3 = k$ schneiden den Quader $0 \le x_1 \le 6$, $0 \le x_2 \le 2$, $0 \le x_3 \le 2$. Bestimme den größten Bereich zwischen zwei Übergangswerten und die Kanten, die $E_k$ dort schneidet.
\teil Bestimme die Spurgeraden $g_k$ der Ebenen $E_k: k x_1 + x_2 + 2x_3 = 4$ in der $x_1x_2$-Ebene. $g_k: \vec{x} = $ \leerfeld
\teil Für alle $k > k_0$ haben $E_k: x_1 + 3x_2 + x_3 = k$ und der Quader $0 \le x_1 \le 4$, $0 \le x_2 \le 2$, $0 \le x_3 \le 6$ keinen gemeinsamen Punkt. Bestimme das kleinstmögliche $k_0$ und begründe. \hfill (Abitur 2023 LK)
\end{teile}
\end{aufgabe}

%% TODO Titel: Ich-kann-Satz fehlt in der Bank (Platzhalter: Kettenname) – Nr. 17
%% TODO Anweisung: ein Satz über der ersten Teilaufgabe fehlt in der Bank – Nr. 17
\begin{aufgabe}{Scharen am Körper – Fehler finden, Begründen}
\begin{teile}
\teil Finde den Fehler. Ab welchem $k_0$ trifft $E_k: x_1 + 2x_2 + x_3 = k$ den Quader $0 \le x_1 \le 3$, $0 \le x_2 \le 2$, $0 \le x_3 \le 4$ für $k > k_0$ nicht mehr? \rechnung{\text{größte Koordinate } &= 4 && \Rightarrow k_0 = 4}
\teil Warum behält die Schnittfigur einer Ebenenschar mit einem Körper zwischen zwei Übergangswerten ihre Gestalt?
\end{teile}
\end{aufgabe}
